\documentstyle[%


righttag%


,11pt%


,amssymb%


,russian%               remove in Eng.


,izvru%                 Set izven% in Eng.


% ,izvru-ks             for brief communications


]{amsart}





% verify indices near Prop. 3





\newcommand{\Aut}{\operatorname{Aut}}


\newcommand{\tts}[1]{}





\theoremstyle{definition}


\theorembodyfont{\rm}


\newtheorem{cornonum}{\hskip\parindent Следствие}


\newtheorem{notenonum}{\hskip\parindent Замечание}


\newtheorem{defnnonum}{\hskip\parindent Определение}


\renewcommand{\thecornonum}{}


\renewcommand{\thenotenonum}{}


\renewcommand{\thedefnnonum}{}





%\hoffset=-15mm%                  For Eng.


\setcounter{page}{13}


\god{2001}


\nomer{2~(465)} %                        Remove in Eng.


%\nomer{Vol.\,45, No.\,2}%               Set in Eng.


%\str{\pageref{begin}--\pageref{end}}%  Set in Eng.


\udk{512.548}





\author{\it А.М.\,Гальмак}


% NB:   ^^ remove in Eng.


\title{Полиадические аналоги теорем Кэли и Биркгофа}





\date{}


%\pagestyle{myheadings}%         Set in Eng.


\pagestyle{plain} %              Remove in Eng.


\begin{document}


%\thispagestyle{plain}%          Set in Eng.


\maketitle


\label{begin}























Исторически $n$-арные группы возникли как обобщение понятия


группы. Поэтому естественным выглядит то, что одной из важных


задач теории $n$-арных групп является нахождение


результатов-аналогов для соответствующих теорем теории групп. При


этом один и тот же бинарный результат может иметь несколько


различных $n$-арных аналогов. В этом можно убедиться на примере


теоремы Кэли, $n$-арные аналоги для которой из [1] и [2] не


совпадают, хотя оба подхода основаны на использовании обычных, т.\,е.


бинарных подстановок. Между тем, при изучении $n$-арных


групп наряду с бинарными подстановками рассматриваются [3], [4] и их


$n$-арные аналоги --- конечные последовательности обычных


подстановок. Поэтому вполне естественной является следующая


задача: получить $n$-арный аналог теоремы Кэли, в которой роль


симметрической группы играла бы $n$-арная группа $n$-арных


подстановок.





Доказанная в данной работе теорема решает поставленную задачу в


самом общем виде. При этом аналог теоремы Кэли из [1] ввиду того,


что при его получении неявно использовались $n$-арные подстановки,


находится среди следствий этой теоремы, каждое из которых само


является $n$-арным аналогом теоремы Кэли.





Пусть $A$ --- $n$-арная группа с $n$-арной операцией $f$. На множестве


$A^{n-1}$ определим $n$-арную операцию $g$ аналогично $n$-арной операции,


введенной Постом для $n$-арных подстановок,


\begin{multline*}


g((a_1',\dotsc,a_{n-1}'),(a_1'',\dotsc,a_{n-1}''),\dotsc,


(a_1^{(n)},\dotsc,a_{n-1}^{(n)}))= \\


%


=(f(a_1',a_2'',\dotsc,a_{n-1}^{(n)},a_1^{(n)}),


f(a_2',\dotsc,a_{n-1}^{(n-2)},a_1^{(n-1)},a_2^{(n)}),\dotsc,


f(a_{n-1}',a_1'',\dotsc,a_{n-1}^{(n)})).


\end{multline*}





Ясно, что декартова степень $A^{n-1}$ вместе с $n$-арной


операцией $g$ является $n$-арной группой.


В $n$-арной группе


$A^{n-1}$ выделим подмножество


$A_0=\{(\,\underbrace{a,\dotsc,a}_{n-1}\,)\mid a\in A\}$.


Так как


$$


g((\,\underbrace{a_1,\dotsc,a_1}_{n-1}\,),\dotsc,


(\,\underbrace{a_n,\dotsc,a_n}_{n-1}\,))=


(f(a_1,\dotsc,a_n),\dotsc,f(a_1,\dotsc,a_n))\in A_0,


$$


то множество $A_0$ замкнуто относительно операции $g$.





\begin{prop}


$A_0$ вместе с $n$-арной операцией $g$


является $n$-арной группой, изоморфной $n$-арной группе $A$.


\end{prop}





\begin{pf}


Определим отображение $\alpha: A\to A_0$ по


правилу $\alpha:a\mapsto (\,\underbrace{a,\dotsc,a}_{n-1}\,)$.


Ясно, что $\alpha$ --- биекция. Так как


\begin{multline*}


\alpha(f(a_1,\dotsc,a_n))=


(\,\underbrace{f(a_1,\dotsc,a_n),\dotsc,f(a_1,\dotsc,a_n)}_{n-1}\,)= \\


%


=g((\,\underbrace{a_1,\dotsc,a_1}_{n-1}\,),\dotsc,


(\,\underbrace{a_n,\dotsc,a_n}_{n-1}\,))=


g(\alpha(a_1),\dotsc,\alpha(a_n)),


\end{multline*}


то $\alpha$ --- изоморфизм $n$-арной группы $A$ на $n$-арную группу


$A_0$.


\end{pf}





Для $n$-арной группы $A$ с $n$-арной операцией $f$ на множестве $A^i$


$(i = 1,\dotsc, n-1)$ определим отношение $\Theta_i$ по правилу


$(a_1,...,a_i)\Theta_i(b_1,...,b_i)$


% return ^^ \dotsc in   ^^ eng.


тогда и только тогда, когда


$f(a_1,...,a_i,x_{i+1},...,x_n)=f(b_1,\dotsc,b_i,x_{i+1},\dotsc,x_n)$


% return ^^ \dotsc in   ^^ eng.


для любых $x_{i+1},\dotsc, x_n \in A$. Отношение $\Theta_i$


является эквивалентностью на $A^i$ ([3], c.\,17; [4], c.\,216--218).


Класс эквивалентности, определяемый элементом $(a_1,\dotsc,a_i)$,


обозначим через $\Theta_i[a_1,\dotsc,a_i]$, положим также


$A_i = A^i/\Theta_i$. Можно показать, что


\begin{equation}


A_i=\{\Theta_i[a_1,\dotsc,a_{i-1},c]\mid c\in A\} =


\{\Theta_i [c,a_1,\dotsc,a_{i-1}] \mid c\in A\},


\end{equation}


где $a_1,\dotsc, a_{i-1}$ --- фиксированные элементы из $A$.





Пусть $\sigma = (1,2,\dotsc,n-1)$ --- циклическая подстановка. Для


произвольного элемента $c\in A$ и любого $i = 1,\dotsc,n-1$ определим


отображения


$r_{ci}: A_i\to A_{\sigma(i)}$, $l_{ci}: A_i\to A_{\sigma(i)}$


по формулам


\begin{gather*}


r_{ci}(\Theta_i [a_1,\dotsc,a_i]) = \Theta_{i+1}


[a_1,\dotsc,a_i,c],\quad i = 1,\dotsc, n-2, \\


%


r_{c(n-1)}(\Theta_{n-1}[a_1,\dotsc,a_{n-1}]) = \Theta_1


[f(a_1,\dotsc,a_{n-1},c)], \\


%


l_{ci}(\Theta_i [a_1,\dotsc,a_i]) = \Theta_{i+1} [c,a_1,\dotsc,a_i],


\quad i = 1,\dotsc, n-2, \\


%


l_{c(n-1)}(\Theta_{n-1} [a_1,\dotsc,a_{n-1}]) = \Theta_1


[f(c,a_1,\dotsc,a_{n-1})].


\end{gather*}


Используя (1), можно показать, что $r_{ci}$ и $l_{ci}$ --- биекции.





Положим


\begin{alignat*}{2}


r_{c_1,\dotsc,c_{n-1}}&=(r_{c_11},\dotsc,r_{c_{n-1}(n-1)}),&\quad


R(A)&=\{r_{c_1,\dotsc,c_{n-1}}\mid c_1,\dotsc,c_{n-1}\in A\}, \\


%


l_{c_1,\dotsc,c_{n-1}}&=(l_{c_11},\dotsc,l_{c_{n-1}(n-1)}),&\quad


L(A)&=\{l_{c_1,\dotsc,c_{n-1}}\mid c_1,\dotsc,c_{n-1}\in A\}.


\end{alignat*}





Ясно, что $r_{c_1,\dotsc,c_{n-1}}$  и


$l_{c_1,\dotsc,c_{n-1}}$ для любых $c_1,\dotsc,c_{n-1}\in A$ являются


$n$-арными подстановками. Поэтому $R(A)$ и $L(A)$ являются


подмножествами множества $S_{A_1,\dotsc, A_{n-1}}$ всех $n$-арных подстановок


последовательности $(A_1,\dotsc, A_{n-1})$.





На множестве $R(A)$ определим $n$-арную операцию $h$ по правилу


$$


h(r_{c_1',\dotsc,c_{n-1}'},r_{c_1'',\dotsc,c_{n-1}''},\dotsc,r_{c_1^{(n)}


,\dotsc,c_{n-1}^{(n)}})=


r_{f(c_1',c_2'',\dotsc,c_{n-1}^{(n-1)},c_1^{(n)}),


f(c_2',\dotsc,c_{n-1}^{(n-2)},c_1^{(n-1)},c_2^{(n)}),\dotsc,


f(c_{n-1}',c_1'',\dotsc,c_{n-1}^{(n)})},


$$


где


$$


r_{c_1^{(j)},\dotsc,c_{n-1}^{(j)}}=


(r_{c_1^{(j)}1},\dotsc,r_{c_{n-1}^{(j)}(n-1)}),\quad


j = 1,\dotsc, n.


$$





\begin{prop}


Множество $R(A)$ вместе с $n$-арной


операцией $h$ является $n$-арной группой.


\end{prop}





{\bf Доказательство.} Легко заметить, что $n$-арная операция $h$


совпадает с ассоциативной $n$-арной операцией, определенной Постом


на множестве $S_{A_1,\dotsc,A_{n-1}}$ всех $n$-арных подстановок


последовательности $(A_1,\dotsc, A_{n-1})$ произвольных множеств


одинаковой мощности.





Полагая $c_1'=x_1$ --- решение уравнения $f(x_1,c_2'',\dotsc,c_{n-1}^{(n-1)},


c_1^{(n)})=c_1,\dotsc,c_{n-1}'=x_{n-1}$ --- решение


уравнения $f(x_{n-1},c_1'',\dotsc,c_{n-1}^{(n)})=c_{n-1}$,


убеждаемся, что $u=r_{c_1',\dotsc,c_{n-1}'}$  является решением уравнения


$$


h(u,r_{c_1'',\dotsc,c_{n-1}''},\dotsc,


r_{c_1^{(n)},\dotsc,c_{n-1}^{(n)}})=


r_{c_1,\dotsc,c_{n-1}}.


$$


Полагая $c_1^{(n)}=y_1$ --- решение уравнения


$f(c_1',c_2'',\dotsc,c_{n-1}^{(n-1)},y_1)=c_1,\dotsc,


c_{n-1}^{(n)}=y_{n-1}$ --- решение уравнения


$f(c_{n-1}',c_1'',\dotsc,c_{n-2}^{(n-1)},y_{n-1})=c_{n-1}$,


убеждаемся, что $v=r_{c_1^{(n)},\dotsc,c_{n-1}^{(n)}}$ является


решением уравнения


$$


h(r_{c_1',\dotsc,c_{n-1}'},\dotsc,


r_{c_1^{(n-1)},\dotsc,c_{n-1}^{(n-1)}},v)=r_{c_1,\dotsc,c_{n-1}}.


\quad\square


$$





На $L(A)$ $n$-арная операция $h$ определяется так же, как на $R(A)$,


и аналогично доказывается, что множество $L(A)$ вместе с $n$-арной


операцией $h$ является $n$-арной группой.





\begin{defnnonum}


Элементы множества $R(A)$ назовем правыми


$n$-арными сдвигами $n$-арной группы $A$. Элементы множества


$L(A)$ назовем левыми $n$-арными сдвигами $n$-арной группы $A$.


\end{defnnonum}





Среди всех $n$-арных сдвигов выделим правые $n$-арные сдвиги вида


$(r_{c1},\dotsc,r_{c(n-1)}){=}


r_{\underbrace{\scriptstyle{c,\dotsc,c}}_{n-1}}{=}r_c$


и левые $n$-арные сдвиги вида


$(l_{c1},\dotsc,l_{c(n-1)})=


l_{\underbrace{\scriptstyle{c,\dotsc,c}}_{n-1}}=l_c.$


Положим $R_0(A)=\{r_c\mid c\in A\}$, $L_0(A) = \{l_c\in c\in A\}$.





\begin{prop}


Множества $R_0(A)$ и $L_0(A)$ являются


$n$-арными подгруппами соответственно $n$-арных групп $R(A)$ и


$L(A)$.


\end{prop}





\begin{pf}


Согласно определению $n$-арной


операции $h$ для любых $r_{c_1},\dotsc,r_{c_n}\in R_0(A)$ будем иметь


$$


h(r_{c_1},\dotsc,r_{c_n})=


h(r_{\underbrace{\scriptstyle{c_1,\dotsc,c_1}}_{n-1}},\dotsc,


r_{\underbrace{\scriptstyle{c_n,\dotsc,c_n}}_{n-1}}\,)=


r_{\underbrace{\scriptstyle{f(c_1,\dotsc,c_n),\dotsc,


f(c_1,\dotsc,c_n)}}_{n-1}}=r_{f(c_1,\dotsc,c_n)},


$$


т.\,е. $h(r_{c_1},\dotsc,r_{c_n})=r_{f(c_1,\dotsc,c_n)}$.


Следовательно, множество $R_0(A)$ замкнуто относительно


ассоциативной $n$-арной операции $h$.





Разрешимость уравнений


$$


h(u,r_{c_2},\dotsc,r_{c_n})=r_c,\quad


h(r_{c_1},\dotsc,r_{c_{n-1}},v)=r_c


$$


является следствием разрешимости в $n$-арной группе $A$ уравнений


$$


f(x,c_2,\dotsc,c_n) = c,\quad f(c_1,\dotsc,c_{n-1},y) = c.


$$


Мы показали, что $R_0(A)$ --- $n$-арная подгруппа $n$-арной группы


$R(A)$. Для $L_0(A)$ доказательство проводится аналогично.


\end{pf}





\begin{notenonum}


Ясно, что при $n = 2$ группы $R(A)$ и $R_0(A)$,


соответственно группы $L(A)$ и $L_0(A)$, совпадают.


\end{notenonum}





Пост заметил ([4], с.\,248), что если равномощные множества


$A_1,\dotsc,A_{n-1}$ попарно не пересекаются, то


всякой $n$-арной подстановке  $t=(t_1,\dotsc,t_{n-1})$


последовательности $(A_1,\dotsc, A_{n-1})$ можно поставить в


соответствие бинарную подстановку $\tau$ множества


$\bigcup\limits_{i=1}^{n-1}A_i$, которая действует на $A_i$ так же, как $t_i$.


Обозначив через $\wt S_{A_1,\dotsc,A_{n-1}}$


множество всех таких подстановок и определив $n$-арную операцию


$\wt h(\tau_1,\tau_2,\dotsc,\tau_n)=\tau_1\tau_2\dotsc\tau_n$ для любых


$\tau_1,\tau_2,\dotsc, \tau_n\in \wt S_{A_1,\dotsc,A_{n-1}}$,


можно показать, что $\wt S_{A_1,\dotsc,A_{n-1}}$


вместе с $n$-арной операцией $\wt h$ является


$n$-арной группой, изоморфной $n$-арной группе $S_{A_1,\dotsc,A_{n-1}}$.





В $n$-арной группе $\wt S_{A_1,\dotsc,A_{n-1}}$


естественно выделяются $n$-арные подгруппы $\wt R(A)$, $\wt L(A)$,


$\wt R_0(A)$ и $\wt L_0(A)$, которые изоморфны


соответственно $n$-арным группам $R(A)$, $L(A)$, $R_0(A)$ и $L_0(A)$.





\begin{thm}


Для всякой $n$-арной группы $\langle A,f\rangle $ существует


изоморфизм $n$-арной группы $\langle A^{n-1},g\rangle $ на


$n$-арную группу $\langle R(A),h\rangle $.


\end{thm}





\begin{pf}


Определим отображение


$\beta:A^{n-1}\to R(A)$ по правилу


$$


\beta: (c_1, \dotsc, c_{n-1})\mapsto r_{c_1,\dotsc,c_{n-1}}.


$$


Ясно, что $\beta$ --- сюръекция. Предположим, что


$r_{c_1,\dotsc,c_{n-1}}=r_{b_1,\dotsc,b_{n-1}}$.


Это означает, что $r_{c_ii}=r_{b_ii}$ для любого $i=1,\dotsc, n-1$, т.\,е.


$$


r_{c_ii}(\Theta_i[a_1,\dotsc,a_i])=


r_{b_ii}(\Theta_i[a_1,\dotsc,a_i])


$$


для любого класса эквивалентности $\Theta_i[a_1,\dotsc, a_i]\in A_i$.


Из последнего равенства имеем


\begin{alignat*}{2}


\Theta_{i+1}[a_1,\dotsc,a_i,c_i]&=\Theta_{i+1}[a_1,\dotsc,a_i,b_i]&\quad


\text{при}\ \; i &= 1,\dotsc, n-2; \\


%


\Theta_1[f(a_1,\dotsc, a_{n-1},c_{n-1})] &= \Theta_1[f(a_1,\dotsc, a_{n-1},


b_{n-1})]&\quad \text{при}\ \; i &= n-1,


\end{alignat*}


откуда вытекает


\begin{gather*}


(a_1,\dotsc, a_i,с_i)\Theta_{i+1}(a_1,\dotsc, a_i,b_i)\quad \text{при}


\ \; i = 1,\dotsc, n-2; \\


%


f(a_1,\dotsc, a_{n-1},c_{n-1}) = f(a_1,\dotsc, a_{n-1},b_{n-1})\quad


\text{при}\ \; i =n-1.


\end{gather*}


Учитывая определение отношения $\Theta_i$, окончательно получаем


$c_i = b_i$ для любого $i = 1,\dotsc, n-1$, т.\,е. $(c_1,\dotsc,


c_{n-1}) = (b_1,\dotsc, b_{n-1})$. Мы показали, что $\beta$ ---


инъекция, а значит, и биекция.





Так как


\begin{multline*}


\beta(g((c_1',\dotsc,c_{n-1}'),\dotsc,(c_1^{(n)},\dotsc,c_{n-1}^{(n)})))=


\beta(f(c_1',\dotsc,c_{n-1}^{(n-1)},c_1^{(n)}),\dotsc,


f(c_{n-1}',c_1'',\dotsc,c_{n-1}^{(n)}))= \\


%


=r_{f(c_1',\dotsc,c_{n-1}^{(n-1)},c_1^{(n)}),\dotsc,


f(c_{n-1}',c_1'',\dotsc,c_{n-1}^{(n)})}=


h(r_{c_1',\dotsc,c_{n-1}'},\dotsc,


r_{c_1^{(n)},\dotsc,c_{n-1}^{(n)}})= \\


%


=h(\beta(c_1',\dotsc,c_{n-1}'),\dotsc,


\beta(c_1^{(n)},\dotsc,c_{n-1}^{(n)})),


\end{multline*}


то $\beta$ --- изоморфизм.


\end{pf}





\begin{notenonum} Теорема 1, как и другие $n$-арные аналоги


групповых результатов, может быть доказана при помощи теоремы


Поста о смежных классах~[4] с использованием соответствующего


бинарного прототипа, в данном случае --- теоремы Кэли для групп.


\end{notenonum}





Ясно, что $\beta(A_0) = R_0(A)$, т.\,е. сужение $\beta$ на $A_0$


является изоморфизмом $n$-арных групп


$\langle A_0,g\rangle $ и $\langle R_0(A),h\rangle$.





\begin{cor}


Всякая $n$-арная группа $\langle A,f\rangle $ изоморфна


$n$-арной группе $\langle R_0(A),h\rangle $.


\end{cor}





Искомый изоморфизм определяется произведением $\alpha\beta$, где


$\alpha$ --- изоморфизм предложения~1, $\beta$ --- изоморфизм из


теоремы~1.





\begin{cor}


Для всякой $n$-арной группы $\langle A, f\rangle $


существует изоморфизм $n$-арной группы $\langle A^{n-1},g\rangle $ на


$n$-арную группу $\langle \wt R(A),\wt h\rangle$.


\end{cor}





Искомый изоморфизм равен произведению $\beta\gamma$, где $\gamma$


--- изоморфизм $n$-арной группы $\langle R(A),h\rangle $ на $n$-арную группу


$\langle \wt R(A),\wt h\rangle $.





\begin{cor}[{\series{m}\shape{n}\selectfont [1]}] Всякая $n$-арная группа


$\langle A,f\rangle $ изоморфна $n$-арной группе


$\langle \wt R_0(A),\wt h\rangle $.


\end{cor}





Искомый изоморфизм равен произведению $\alpha\beta\gamma$.





Теорема~1 и каждое из следствий 1--3 являются аналогами


теоремы Кэли для полиадических групп. Нетрудно заметить, что при


$n=2$ эта теорема и все ее следствия совпадают с теоремой Кэли.





Напомним, что $n$-арным автоморфизмом последовательности


$\{A_1,\dotsc,A_{n-1}\}$ однотипных универсальных алгебр называется~[5]


последовательность $f = \{f_1,\dotsc, f_{n-1}\}$ изоморфизмов


$$


A_1@>{f_1}>> A_2 @>{f_2}>> \dotsb


@>{f_{n-2}}>> A_{n-1} @>{f_{n-1}}>> A_1.


$$


Множество всех $n$-арных автоморфизмов последовательности


$\{A_1,\dotsc,A_{n-1}\}$ обозначим через $\Aut(A_1,\dotsc, A_{n-1})$.





Следующая теорема и следствия из нее являются $n$-арными


аналогами известной теоремы Биркгофа для групп.





\begin{thm}


Для всякой $n$-арной группы $\langle A,f\rangle $


существует изоморфизм $n$-арной группы $\langle A^{n-1},g\rangle $ на


$n$-арную группу всех $n$-арных автоморфизмов некоторой


последовательности универсальных алгебр.


\end{thm}





\begin{pf}


Для любого $b\in A$ определим преобразование


$\varphi_b :\bigcup\limits_{i=1}^{n-1}A_i\to\bigcup\limits_{i=1}^{n-1}A_i$


по правилу


$$


\varphi_b : \Theta_i[a_1,\dotsc, a_i]\mapsto\Theta_i[f(b, b_1,\dotsc, b_{n-2},


a_1), a_2,\dotsc, a_i],


$$


где $b_1,\dotsc,b_{n-2}$ --- фиксированные элементы из $A$.


Положим $\Omega = \{\varphi_b \mid b\in A\}$ и


рассмотрим последовательность универсальных алгебр


\begin{equation}


\langle A_1, \Omega \rangle ,\dotsc,\langle A_{n-1}, \Omega\rangle .


\end{equation}


Так как


\begin{multline*}


r_{c_ii}(\varphi_b (\Theta_i[a_1,\dotsc,a_i])) =


r_{c_ii}(\Theta_i[f(b, b_1,\dotsc, b_{n-2}, a_1), a_2,\dotsc, a_i]) = \\


%


= \Theta_{i+1}[f(b, b_1,\dotsc, b_{n-2}, a_1), a_2,\dotsc, a_i, c_i] =


\varphi_b(\Theta_{i+1}[a_1,\dotsc, a_i, c_i]) =


\varphi_b(r_{c_ii}(\Theta_i[a_1,\dotsc, a_i]))


\end{multline*}


для любой операции $\varphi_b\in \Omega$, то $r_{c_ii}$ ---


изоморфизм алгебры $\langle  A_i, \Omega \rangle $ на алгебру


$\langle  A_{i+1}, \Omega \rangle $


$(i = 1, \dotsc, n-2)$.





Аналогично доказывается, что $r_{c_{(n-1)}(n-1)}$ --- изоморфизм алгебры


${\langle A_{n-1}, \Omega \rangle }$ на алгебру


$\langle A_1, \Omega \rangle $.


Следовательно, $r_{c_1,\dotsc,c_{n-1}}$ --- $n$-арный


автоморфизм последовательности (2). Мы доказали включение


\begin{equation}


R(A) \subseteq \Aut(A_1,\dotsc, A_{n-1}).


\end{equation}





Пусть теперь $\delta = (\delta_1,\dotsc, \delta_{n-1})\in


\Aut(A_1,\dotsc, A_{n-1})$, т.\,е. все $\delta_i$ ---


биекции и верно


\begin{equation}


\delta_i(\varphi_b (\Theta_i[a_1,\dotsc, a_i])) =


\varphi_b(\delta_i(\Theta_i[a_1,\dotsc, a_i])),\quad


i = 1,\dotsc, n-1,


\end{equation}


для любой операции $\varphi_b\in \Omega$ и любого класса


$\Theta_i[a_1,\dotsc, a_i]\in A_i$. Последнее равенство можно переписать


следующим образом:


\begin{equation}


\delta_i(\Theta_i[f(b, b_1,\dotsc,b_{n-2}, a_1),


a_2,\dotsc,a_i]) =


\Theta_{i+1}[f(b, b_1,\dotsc, b_{n-2},d_1),


d_2,\dotsc,d_{i+1}]


\end{equation}


при $i = 1,\dotsc,n-2$, где $\delta_i(\Theta_i[a_1,\dotsc, a_i]) =


\Theta_{i+1}[d_1,\dotsc, d_{i+1}]$;


\begin{equation}


\delta_{n-1}(\Theta_{n-1}[f(b, b_1,\dotsc, b_{n-2}, a_1), a_2,\dotsc,


a_{n-1}]) = \Theta_1[f(b, b_1,\dotsc, b_{n-2}, d)]


\end{equation}


при $i = n-1$, где $\delta_{n-1}(\Theta_{n-1}[a_1,\dotsc, a_{n-1}]) =


\Theta_1[d]$.





Так как (4) справедливо для любого класса $\Theta_i[a_1,\dotsc, a_i]\in A_i$,


то элемент $a_1$ можно выбрать так, что $(b_1,\dotsc, b_{n-2}, a_1)$


--- нейтральная последовательность. С учетом этого (5) и~(6)


примут соответственно вид


\begin{align}


\delta_i(\Theta_i[b, a_2,\dotsc, a_i]) &=


\Theta_{i+1}[f(b, b_1,\dotsc, b_{n-2}, d_1), d_2,\dotsc, d_{i+1}], \\


%


\delta_{n-1}(\Theta_{n-1}[b, a_2,\dotsc, a_{n-1}]) &=


\Theta_1[f(b, b_1,\dotsc, b_{n-2}, d)].


\end{align}


В $n$-арной группе всегда существуют элементы $c_i$ такие, что


$$


(d_1,\dotsc, d_{i+1})\Theta_i (a_1,\dotsc, a_i, c_i),


$$


и элемент $c_{n-1}$ такой, что $d = f(a_1,\dotsc, a_{n-1}, c_{n-1})$.


Теперь~(7) и~(8) перепишутся в виде


\begin{align*}


\delta_i(\Theta_i[b, a_2,\dotsc, a_i]) &=


\Theta_{i+1}[f(b, b_1,\dotsc, b_{n-2}, a_1), a_2,\dotsc, a_i, c_i], \\


%


\delta_{n-1}(\Theta_{n-1}[b, a_2,\dotsc, a_{n-1}]) &=


\Theta_1[f(b, b_1,\dotsc, b_{n-2}, a_1), a_2,\dotsc, a_{n-1}, c_{n-1}].


\end{align*}


Учитывая нейтральность последовательности $(b_1,\dotsc, b_{n-2}, a_1)$, из


полученных равенств имеем


\begin{align*}


\delta_i(\Theta_i[b, a_2,\dotsc, a_i]) &=


\Theta_{i+1}[b, a_2,\dotsc, a_i, c_i],\quad i = 1,\dotsc, n-2, \\


%


\delta_{n-1}(\Theta_{n-1}[b, a_2,\dotsc, a_{n-1}]) &=


\Theta_1[f(b, a_2,\dotsc, a_{n-1}, c_{n-1})].


\end{align*}


Первое из этих равенств справедливо для всех элементов


множества $A_i$ $(i=1,\dotsc,n-2)$, а второе справедливо для всех элементов


множества $A_{n-1}$. Следовательно, отображение $\delta_i$ совпадает с


отображением  $r_{c_ii}$ $(i=1,\dotsc,n-1)$, а $\delta = (\delta_1,\dotsc,


\delta_{n-1})$ является правым $n$-арным сдвигом.


В силу произвольного выбора $\delta\in \Aut(A_1,\dotsc, A_{n-1})$


доказано включение


$$


\Aut(A_1,\dotsc, A_{n-1}) \subseteq R(A).


$$


Oтсюда и из (3) получаем


$$


R(A) = \Aut(A_1,\dotsc, A_{n-1}).


$$


Применяя теорему~1, получаем изоморфизм $\beta$ $n$-арной группы


$\langle A^{n-1},g\rangle$


на $n$-арную группу $\langle \Aut (A_1,\dotsc, A_{n-1}), h \rangle $.


\end{pf}





Так как $R_0(A)\subseteq R(A) =\Aut(A_1,\dotsc, A_{n-1})$ и


согласно следствию~1 существует изоморфизм $n$-арных групп $\langle  A,


f \rangle $ и $\langle R_0, h \rangle $, то справедливо





\begin{cor}


Всякая $n$-арная группа изоморфна


$n$-арной группе $n$-арных автоморфизмов некоторой


последовательности универсальных алгебр.


\end{cor}





\begin{cor}


Для всякой $n$-арной группы $\langle  A, f \rangle $


существует изоморфизм $n$-арной группы $\langle A^{n-1}, g \rangle $ на


$n$-арную группу автоморфизмов некоторой универсальной алгебры.


\end{cor}





\begin{pf}


Согласно следствию~2 существует


изоморфизм $\mu =\beta\gamma$ $n$-арной группы $\langle A^{n-1}, g \rangle$


на $n$-арную группу $\langle \wt R(A),\wt h\rangle$, где


\begin{alignat*}{2}


\beta: A^{n-1} &\to R(A),&\quad \beta: (c_1,\dotsc,c_{n-1})&\mapsto


r_{c_1,\dotsc,c_{n-1}}; \\


%


\gamma: R(A)&\to\wt R(A), &\quad


\gamma:r_{c_1,\dotsc,c_{n-1}}&\mapsto r,


\end{alignat*}


причем $r$ --- биекция множества $\bigcup\limits_{i=1}^{n-1}A_i$,


которая действует на $A_i$ так же, как $r_{c_ii}$ $(i=1,\dotsc, n-1)$.





Рассмотрим универсальную алгебру


$\Big\langle \bigcup\limits_{i=1}^{n-1}A_i,\Omega\Big\rangle $. При


доказательстве теоремы~2 было установлено, что все $r_{c_ii}$ ---


изоморфизмы. Поэтому


$$


r(\varphi_b (\Theta_i[a_1,\dotsc, a_i])) =


r_{c_ii}(\varphi_b (\Theta_i[a_1,\dotsc, a_i])) =


\varphi_b(r_{c_ii}(\Theta_i[a_1,\dotsc, a_i])) =


\varphi_b (r(\Theta_i[a_1,\dotsc, a_i]))


$$


для любого $\Theta_i[a_1,\dotsc, a_i]\in \bigcup\limits_{i=1}^{n-1}A_i$


и любой операции $\varphi_b\in \Omega$, т.\,е. $r$ ---


автоморфизм алгебры


$\Big\langle \bigcup\limits_{i=1}^{n-1}A_i,\Omega\Big\rangle$.


\end{pf}





Следствию 3 из теоремы 1 соответствует





\begin{cor}


Всякая $n$-арная группа изоморфна $n$-арной группе


автоморфизмов некоторой универсальной алгебры.


\end{cor}





\begin{thebibliography}{9}





\bibitem{1}


Sioson F.M. {\em On free Abelian $m$-groups}. I // Proc. Japan Acad. --


1967. -- V.\,43. -- P.\,876--879.


\bibitem{2}


Гальмак А.М. {\em Трансляции $n$-арных групп} // ДАН БССР. -- 1986. --


Т.\,30. -- \No\,8. -- С.\,677--680.


\bibitem{3}


Русаков С.А. {\em Алгебраические $n$-арные системы}. -- Минск: Навука i


тэхнiка, 1992. -- 264~с.


\bibitem{4}


Post E.L. {\em Polyadic groups} // Trans. Amer. Math. Soc. -- 1940. --


V.\,48. -- \No\,2. -- P.\,208--350.


\bibitem{5}


Гальмак А.М. {\em $n$-aрные морфизмы алгебраических систем} // Тез. докл.


Международн. матем. конф. -- Минск, 1993. -- С.\,11--12.





\end{thebibliography}





\vskip 0.5cm





\postup{\it Могилевский технологический институт}{\qquad \it Поступили}


\postup{}{\it первый вариант $21.07.1997$}


\postup{}{\it окончательный вариант $19.04.2000$}


\label{end}


\end{document}











